\documentclass[10pt, letterpaper]{article}

\usepackage{../../../../template/template}

\setDocTitle{Verbale interno 16/03/2026}
\setDocVersion{-}
\setDocData{16/03/2026}
\setDocIndex{vi\_2026\_03\_16}
\setDocState{1}
\setDocRecipient{1}

\begin{document}

\makeTitlePage

% \newpage
% \addChangelog{0.1.0}{11/03/2026}{F. Pasqual}{C. Libralato}{\noOne}{\firstDraft}
% \makeChangelog

\newpage

\tableofcontents

\newpage

% -------------------------------------------------
% Informazioni riunione
% -------------------------------------------------
\makeMeetingInfoTable{16/03/2026}{14:30}{16:30}{via \textit{Discord$_G$}}

\addParticipant{Valeria}{Baleanu}{Programmatore}{P}
\addParticipant{Leonardo}{Pellizzon}{Programmatore}{P}
\addParticipant{Luca}{Granziero}{Programmatore}{P}
\addParticipant{Francesco}{Pasqual}{Responsabile}{P}
\addParticipant{Giuseppe}{De Fina}{Programmatore}{P}
\addParticipant{Christian}{Libralato}{Verificatore}{P}
\addParticipant{Filippo}{Venzo}{Amministratore}{P}
\makeMeetingParticipantsTable


\addPrecedentTodo{vi\_2026\_03\_09.a1}{\#161}{F. Pasqual, C. Libralato}{16/03/2026}
\addPrecedentTodo{vi\_2026\_03\_09.a2}{\#163}{F. Pasqual, C. Libralato}{11/03/2026}
\addPrecedentTodo{vi\_2026\_03\_09.a3}{\#164}{F. Venzo}{22/03/2026}
\addPrecedentTodo{vi\_2026\_03\_09.a4}{\#166}{F. Pasqual, C. Libralato}{11/03/2026}
\addPrecedentTodo{vi\_2026\_03\_09.a5}{\#167}{V. Baleanu, L. Pellizzon, L. Granziero, G. De Fina}{18/03/2026}
\addPrecedentTodo{vi\_2026\_03\_09.a6}{\#165}{F. Venzo}{22/03/2026}
\addPrecedentTodo{vi\_2026\_03\_09.a7}{\#169}{F. Pasqual}{16/03/2026}
\addPrecedentTodo{vi\_2026\_03\_09.a8}{\#172}{V. Baleanu, L. Pellizzon, L. Granziero, G. De Fina}{16/03/2026}
\addPrecedentTodo{vi\_2026\_03\_09.a9}{\#174}{V. Baleanu, L. Pellizzon, L. Granziero, G. De Fina}{16/03/2026}
\addPrecedentTodo{vi\_2026\_03\_09.a10}{\#173}{F. Pasqual}{16/03/2026}
\makePrecedentTodoTable
% -------------------------------------------------
% Ordine del giorno
% -------------------------------------------------
% -------------------------------------------------
% Ordine del giorno
% -------------------------------------------------
\section{Ordine del giorno}
\begin{itemize}
    \item Implementazione autenticazione \textit{JWT$_G$} e gestione \textit{refresh token};
    \item organizzazione progetto: struttura moduli, cartelle e convenzioni;
    \item scelta delle versioni \textit{Node.js$_G$} e \textit{NestJS$_G$} per il progetto;
    \item gestione del \textit{workflow} e del \textit{branching$_G$}: divisione \textit{frontend}/\textit{backend};
    \item divisione dei compiti e scadenze provvisorie sui moduli utenti, allarmi e \textit{analytics}.
\end{itemize}

% -------------------------------------------------
% Approfondimento
% -------------------------------------------------
\section{Approfondimento}

\subsection*{Implementazione autenticazione \textit{JWT$_G$} e gestione \textit{refresh token}}
Il team ha discusso l'implementazione del sistema di \textit{login} e \textit{JWT} nel \textit{backend}.
Si è valutata la necessità di includere o meno il \textit{refresh token} e i metodi associati.
A corredo, è stato condiviso un documento di riferimento per unificare le regole di validazione
di \textit{username} e altri dati tra le due componenti.

\subsection*{Organizzazione progetto: struttura moduli, cartelle e \textit{convention}}
Si è affrontata la definizione della struttura del progetto \textit{backend} in \textit{NestJS} applicando i principi
dell'architettura esagonale. È stata proposta e approvata la suddivisione delle cartelle in \textit{application}
(che include porte e \textit{use case$_G$}), \textit{domain} (\textit{core}) e \textit{infrastructure}
(\textit{repository}, \textit{entity}, \textit{DTO$_G$}).
Si è convenuto di separare i \textit{DTO} nello strato di infrastruttura e di mantenere la convenzione dei nomi in
minuscolo.

\subsection*{Scelta delle versioni \textit{Node.js}, \textit{NestJS} e \textit{Angular$_G$} per il progetto}
A seguito di una discussione tecnica sulle versioni degli ambienti da utilizzare, è stato deciso di adottare
\textit{Node.js} versione 24.14 LTS, in quanto garantisce supporto a lungo termine e maggiore stabilità fino al 2028.
Per quanto riguarda il \textit{framework backend}, \textit{NestJS} verrà utilizzato alla versione 11 stabile.
Per quanto riguarda il \textit{framework frontend}, \textit{Angular} verrà utilizzato alla versione 21 stabile.
È stata evidenziata l'assoluta necessità di uniformare dipendenze e versioni tra tutti i membri per evitare incompatibilità.
Infine, si è discusso dei futuri strumenti di
\textit{CI/CD$_G$} e \textit{caching} per le \textit{build}, valutando soluzioni \textit{monorepo$_G$} come \textit{NX} o
\textit{Turbo Repo}.


\subsection*{Gestione \textit{workflow} e \textit{branching}}
Per organizzare il processo di sviluppo, si è deciso di dividere la \textit{codebase} in due cartelle: \textit{backend} e \textit{frontend}.
Allo scopo di mantenere l'indipendenza delle due macro-aree, è stato deciso di separare i rami di sviluppo creando due \textit{branch} distinti: \textit{backend} e \textit{frontend},
i quali fungeranno da \textit{branch} base.
L'approccio lavorativo imporrà di sviluppare le funzionalità in appositi \textit{feature$_G$ branch} derivati dai rispettivi
\textit{branch} base, per poi effettuare il \textit{merge$_G$}
unicamente in seguito al superamento dei test.

\subsection*{Divisione compiti, scadenze e questioni in sospeso}
Sono state assegnate le responsabilità di sviluppo iniziali: L. Pellizzon si occuperà del modulo utenti,
L. Granziero e F. Venzo dell'autenticazione verso le \textit{API Vimar} e della gestione degli allarmi, mentre V. Baleanu e G. De Fina prenderanno
in carico gli \textit{analytics} lato \textit{frontend}. L'obiettivo primario è produrre una prima bozza di codice
funzionante entro la fine della settimana, procedendo per rilasci progressivi.
Restano ancora in sospeso alcune questioni che necessitano di chiarimenti, attesi da un colloquio con il prof. Cardin.
Infine, Responsabile e Amministratore dovranno verificare tecnicamente la fattibilità e l'utilità dell'integrazione di \textit{NX} o \textit{Turbo Repo} prima della sua adozione
definitiva.

% -------------------------------------------------
% Decisioni
% -------------------------------------------------
\addDecision{Implementare l'autenticazione con login, \textit{access token JWT e refresh token}}
\addDecision{Adottare un'architettura esagonale nel \textit{backend} con suddivisione in \textit{application}, \textit{domain} e \textit{infrastructure}, mantenendo i \textit{DTO} nello strato \textit{infrastructure}}
\addDecision{Utilizzare \textit{Node.js} versione 24.14 LTS}
\addDecision{Utilizzare \textit{NestJS} versione 11 stabile per il \textit{backend}}
\addDecision{Adottare \textit{Angular} versione 21 stabile per il \textit{frontend}}
\addDecision{Adottare branch base separati per \textit{frontend} e \textit{backend}}
\makeDecisionTable

% -------------------------------------------------
% Azioni
% -------------------------------------------------
\addTodo{\#}{Implementare gestione utenti lato \textit{backend}}{L. Pellizzon}{22/03/2026}
\addTodo{\#}{Implementare autenticazione alle \textit{API Vimar} e sviluppare modulo allarmi}{L. Granziero, F. Venzo}{22/03/2026}
\addTodo{\#}{Implementare modulo \textit{analytics} lato \textit{frontend}}{V. Baleanu, G. De Fina}{22/03/2026}
\addTodo{\#}{Implementare \textit{feature} di ricezione notifiche lato \textit{frontend}}{G. De Fina}{22/03/2026}
\addTodo{\#206}{Documentare in modo definitivo le versioni dei \textit{framework} e dell'ambiente \textit{Node.js} da utilizzare}{F. Pasqual}{22/03/2026}
\addTodo{\#180}{Testare l'integrazione di \textit{NX} o \textit{Turbo Repo}}{F. Pasqual, F. Venzo}{20/03/2026}
\addTodo{\#207}{Concordare un colloquio con il prof. Cardin per chiarire le questioni aperte sul \textit{design$_G$}}{F. Pasqual}{20/03/2026}
\addTodo{\#208}{Stesura verbale della riunione}{F. Pasqual}{20/03/2026}
\addTodo{\#209}{Aggiornamento glossario con termini emersi nella riunione}{F. Pasqual}{20/03/2026}
\makeTodoTable

\end{document}